\documentclass[12pt, a4paper]{article}

\usepackage{geometry}
\geometry{a4paper, top=2.5cm, bottom=2.5cm, left=2.5cm, right=3cm}
\usepackage{graphicx}
\usepackage{booktabs}
\usepackage{longtable}
\usepackage{array}
\usepackage{multirow}
\usepackage{xcolor}
\usepackage{hyperref}
\usepackage{fancyhdr}
\usepackage{enumitem}
\usepackage{fontspec}
\usepackage{xepersian}
\settextfont[Scale=1.1]{Vazirmatn}
\setlatintextfont{Times New Roman}

\definecolor{primaryorange}{RGB}{234, 88, 12}
\definecolor{lightgray}{RGB}{245, 245, 245}
\definecolor{darkgray}{RGB}{80, 80, 80}
\definecolor{successgreen}{RGB}{22, 163, 74}
\definecolor{warnred}{RGB}{220, 38, 38}

\hypersetup{
    colorlinks=true,
    linkcolor=primaryorange,
    urlcolor=primaryorange
}

\pagestyle{fancy}
\fancyhf{}
\rhead{پروژه پاتوق — سامانه رزرو رستوران}
\lhead{گزارش اسپرینت سوم}
\rfoot{\thepage}
\renewcommand{\headrulewidth}{0.4pt}

\begin{document}

% \pagenumbering{persian}  % disabled for hyperref compatibility

%% ─── صفحه عنوان ────────────────────────────────────────────────────────────
\begin{titlepage}
\centering
\vspace*{2cm}
{\LARGE\bfseries پروژه پاتوق (پاتوق)\\[0.5em]}
{\Large\bfseries سامانه مدیریت و رزرو کافه و رستوران\\[2em]}
\rule{\linewidth}{0.5mm}\\[1em]
{\huge\bfseries\color{primaryorange} گزارش اسپرینت سوم}\\[0.5em]
{\Large تکمیل، امنیت و رفع باگ}\\[1em]
\rule{\linewidth}{0.5mm}\\[2em]
{\large درس: طراحی نرم‌افزار\\[0.5em]}
{\large مدت اسپرینت: ۲ هفته\\[0.5em]}
{\large تاریخ تهیه: تیر ۱۴۰۵}\\[3em]
\vfill
{\normalsize\color{darkgray}
این گزارش پوشش‌دهی اسپرینت سوم پروژه پاتوق را شرح می‌دهد که
در آن امنیت سیستم تقویت شد، باگ‌های شناسایی‌شده رفع گردیدند
و صفحات ناقص تکمیل شدند.}
\end{titlepage}

\tableofcontents
\newpage

%% ─── ۱. هدف اسپرینت ────────────────────────────────────────────────────────
\section{هدف اسپرینت}

اسپرینت سوم با هدف تکمیل سیستم، بهبود امنیت و رفع تمامی باگ‌های شناسایی‌شده
در اسپرینت‌های قبل آغاز شد. این اسپرینت آخرین اسپرینت پروژه است و در پایان
باید سیستم آماده تحویل نهایی باشد.

\subsection{اهداف اصلی}
\begin{itemize}[rightmargin=0.5cm]
  \item پیاده‌سازی چرخش توکن تازه‌سازی (\lr{Refresh Token Rotation})
  \item اضافه کردن محدودیت نرخ (\lr{Rate Limiting}) برای \lr{Endpoint}های احراز هویت
  \item رفع مشکل \lr{CORS} برای پورت \lr{5174}
  \item استانداردسازی فرمت پاسخ‌های خطا (\lr{RFC 7807})
  \item بهبود امنیت \lr{OTP}
  \item رفع باگ مدیریت خطای ورود در کلاینت
  \item تکمیل صفحه \lr{CompleteProfile}
  \item رفع باگ تقویم شمسی
  \item بهبود امنیت ورود ادمین
  \item تکمیل صفحات پروفایل
\end{itemize}

\newpage

%% ─── ۲. Sprint Backlog ─────────────────────────────────────────────────────
\section{\lr{Sprint Backlog}}

\begin{longtable}{|r|r|r|r|}
\hline
\textbf{آیتم} & \textbf{نوع} & \textbf{مسئول} & \textbf{وضعیت} \\
\hline
\endhead
چرخش توکن تازه‌سازی & امنیت & بک‌اند & انجام شد \\
\hline
محدودیت نرخ احراز هویت & امنیت & بک‌اند & انجام شد \\
\hline
محدودیت نرخ عمومی & امنیت & بک‌اند & انجام شد \\
\hline
رفع مشکل \lr{CORS} پورت ۵۱۷۴ & باگ & بک‌اند & انجام شد \\
\hline
فرمت خطا \lr{RFC 7807} & بهبود & بک‌اند & انجام شد \\
\hline
\lr{OTP} یکبار مصرف & امنیت & بک‌اند & انجام شد \\
\hline
حذف \lr{OTP} از لاگ & امنیت & بک‌اند & انجام شد \\
\hline
باگ ورود: \lr{401 handler} & باگ & فرانت‌اند & انجام شد \\
\hline
پیام‌های خطای فارسی & بهبود & فرانت‌اند & انجام شد \\
\hline
تغییر نام \lr{IsNewUser} & بهبود & بک‌اند & انجام شد \\
\hline
صفحه تکمیل پروفایل & ویژگی & فرانت‌اند & انجام شد \\
\hline
رفع باگ تقویم شمسی & باگ & فرانت‌اند & انجام شد \\
\hline
بهبود امنیت ورود ادمین & امنیت & بک‌اند+فرانت & انجام شد \\
\hline
تکمیل پروفایل مشتری & ویژگی & فرانت‌اند & انجام شد \\
\hline
تکمیل پروفایل مدیر رستوران & ویژگی & فرانت‌اند & انجام شد \\
\hline
مدیریت میز (تکمیل) & ویژگی & فرانت‌اند & انجام شد \\
\hline
تست نهایی یکپارچگی & تست & هر دو & انجام شد \\
\hline
\end{longtable}

\newpage

%% ─── ۳. امنیت: Refresh Token Rotation ──────────────────────────────────────
\section{چرخش توکن تازه‌سازی}

\subsection{مشکل}
در اسپرینت اول، توکن‌های تازه‌سازی پس از استفاده باطل نمی‌شدند. این به معنای آن
بود که یک توکن دزدیده‌شده می‌توانست بارها برای دریافت توکن دسترسی جدید استفاده شود.

\subsection{راه‌حل: Token Rotation}
الگوریتم چرخش توکن تازه‌سازی پیاده‌سازی شد:

\begin{enumerate}[rightmargin=0.5cm]
  \item کاربر توکن تازه‌سازی را به \lr{POST /api/auth/refresh} ارسال می‌کند
  \item سرور توکن قدیمی را در پایگاه داده پیدا می‌کند
  \item اگر توکن قبلاً باطل شده باشد، \textbf{تمام} توکن‌های آن کاربر باطل می‌شوند
    (احتمال سرقت توکن)
  \item توکن قدیمی با \lr{IsRevoked = true} علامت‌گذاری می‌شود
  \item توکن دسترسی جدید + توکن تازه‌سازی جدید صادر می‌شود
\end{enumerate}

\subsection{پیاده‌سازی}
کلاس \texttt{\lr{RefreshTokenCommandHandler}} این منطق را پیاده‌سازی می‌کند.
موجودیت \lr{RefreshToken} در پایگاه داده با فیلدهای \lr{IsRevoked} و \lr{RevokedAt}
نگهداری می‌شود.

\newpage

%% ─── ۴. امنیت: Rate Limiting ────────────────────────────────────────────────
\section{محدودیت نرخ درخواست‌ها}

\subsection{پیاده‌سازی}
محدودیت نرخ با \lr{ASP.NET Core Rate Limiting Middleware} پیاده‌سازی شد.
دو پروفایل تعریف شدند:

\begin{longtable}{|r|r|r|r|}
\hline
\textbf{نام پروفایل} & \textbf{حداکثر درخواست} & \textbf{بازه زمانی} & \textbf{کاربرد} \\
\hline
\endhead
\lr{auth} & ۵ درخواست & ۱ دقیقه & \lr{Endpoint}های احراز هویت \\
\hline
\lr{general} & ۱۰۰ درخواست & ۱ دقیقه & سایر \lr{Endpoint}ها \\
\hline
\end{longtable}

\subsection{رفتار در صورت تجاوز}
پس از تجاوز از حد مجاز، پاسخ \lr{429 Too Many Requests} با پیام فارسی برگردانده می‌شود:

\begin{latin}
\begin{verbatim}
{
  "Success": false,
  "Message": "تعداد درخواست‌های شما از حد مجاز بیشتر شده است.",
  "RetryAfter": "60 seconds"
}
\end{verbatim}
\end{latin}

\subsection{هدف امنیتی}
پروفایل \lr{auth} جلوگیری از:
\begin{itemize}[rightmargin=0.5cm]
  \item حمله \lr{Brute Force} برای \lr{OTP}
  \item حمله \lr{Credential Stuffing} برای ورود با رمز عبور
  \item \lr{OTP Bombing} (ارسال انبوه \lr{OTP} به یک شماره)
\end{itemize}

\newpage

%% ─── ۵. رفع باگ CORS ────────────────────────────────────────────────────────
\section{رفع مشکل CORS برای پورت ۵۱۷۴}

\subsection{مشکل}
\lr{Vite} در برخی شرایط (وقتی پورت ۵۱۷۳ اشغال است) روی پورت ۵۱۷۴ اجرا می‌شود.
پیکربندی \lr{CORS} اولیه فقط \lr{http://localhost:5173} را مجاز می‌دانست.

\subsection{راه‌حل}
پورت ۵۱۷۴ به لیست \lr{AllowedOrigins} در \lr{appsettings.json} اضافه شد:

\begin{latin}
\begin{verbatim}
"AllowedOrigins": [
  "http://localhost:5173",
  "http://localhost:5174"
]
\end{verbatim}
\end{latin}

\newpage

%% ─── ۶. فرمت خطا RFC 7807 ───────────────────────────────────────────────────
\section{استانداردسازی پاسخ خطا با RFC 7807}

\subsection{مشکل}
پاسخ‌های خطا فرمت یکسانی نداشتند. برخی اوقات رشته متن ساده، برخی اوقات
\lr{JSON} بدون ساختار مشخص برگردانده می‌شد.

\subsection{راه‌حل}
\lr{ExceptionHandlingMiddleware} برای رهگیری تمام خطاها و برگرداندن پاسخ استاندارد
\lr{Problem Details (RFC 7807)} پیاده‌سازی شد:

\begin{latin}
\begin{verbatim}
{
  "type": "https://tools.ietf.org/html/rfc7807",
  "title": "Not Found",
  "status": 404,
  "detail": "رزرو با شناسه داده‌شده یافت نشد.",
  "instance": "/api/reservations/abc-123"
}
\end{verbatim}
\end{latin}

\subsection{کدهای وضعیت استاندارد}
\begin{longtable}{|r|r|r|}
\hline
\textbf{کد} & \textbf{نام} & \textbf{کاربرد} \\
\hline
\endhead
\lr{400} & \lr{Bad Request} & خطای اعتبارسنجی (\lr{FluentValidation}) \\
\hline
\lr{401} & \lr{Unauthorized} & توکن نامعتبر یا منقضی \\
\hline
\lr{403} & \lr{Forbidden} & عدم دسترسی بر اساس نقش \\
\hline
\lr{404} & \lr{Not Found} & موجودیت یافت نشد \\
\hline
\lr{409} & \lr{Conflict} & تداخل (مثل رزرو همزمان یک میز) \\
\hline
\lr{429} & \lr{Too Many Requests} & تجاوز از محدودیت نرخ \\
\hline
\lr{500} & \lr{Internal Server Error} & خطای غیرمنتظره \\
\hline
\end{longtable}

\newpage

%% ─── ۷. امنیت OTP ───────────────────────────────────────────────────────────
\section{بهبود امنیت OTP}

\subsection{یکبار مصرف بودن OTP}
در پیاده‌سازی اولیه، کد \lr{OTP} پس از تأیید از \lr{Redis} حذف نمی‌شد.
این باگ رفع شد: پس از تأیید موفق کد، \lr{IDatabase.KeyDelete} فراخوانی می‌شود.

\subsection{حذف OTP از لاگ}
لاگ‌های \lr{Serilog} حاوی کد \lr{OTP} بودند که خطر امنیتی ایجاد می‌کرد.
تمام \lr{log statement}هایی که مقدار \lr{OTP} را ثبت می‌کردند حذف شدند.
به جای آن، یک لاگ با فرمت \texttt{\lr{OTP sent to \{phone[..3]\}***}} جایگزین شد که
شماره موبایل کامل را نشان نمی‌دهد.

\newpage

%% ─── ۸. باگ فرانت‌اند ───────────────────────────────────────────────────────
\section{رفع باگ‌های فرانت‌اند}

\subsection{باگ ۱: مدیریت خطای 401 در client.ts}
\begin{description}[rightmargin=0.5cm]
  \item[مشکل:] هنگامی که توکن دسترسی منقضی می‌شد، \lr{API} پاسخ \lr{401} برمی‌گرداند.
    فرانت‌اند این پاسخ را مدیریت نمی‌کرد و کاربر یک صفحه خطای خام می‌دید.
  \item[راه‌حل:] در \lr{client.ts} (لایه \lr{API}) یک رهگیر (\lr{interceptor}) اضافه شد
    که پاسخ‌های \lr{401} را شناسایی کرده، توکن از \lr{localStorage} حذف می‌شود و
    کاربر به صفحه ورود هدایت می‌شود.
\end{description}

\subsection{باگ ۲: پیام‌های خطای انگلیسی}
\begin{description}[rightmargin=0.5cm]
  \item[مشکل:] پیام‌های خطا از \lr{API} به انگلیسی بودند و مستقیم نمایش داده می‌شدند.
  \item[راه‌حل:] نقشه‌برداری از کدهای خطا به پیام‌های فارسی در \lr{client.ts} اضافه شد.
    همچنین بک‌اند برای برخی خطاهای رایج پیام فارسی مستقیم برمی‌گرداند.
\end{description}

\subsection{باگ ۳: تقویم شمسی در فرم رزرو}
\begin{description}[rightmargin=0.5cm]
  \item[مشکل:] تاریخ انتخاب‌شده از تقویم شمسی به درستی به فرمت \lr{DateOnly} که
    بک‌اند انتظار داشت تبدیل نمی‌شد.
  \item[راه‌حل:] تابع تبدیل تاریخ با استفاده از کتابخانه \lr{date-fns} اصلاح شد.
    تاریخ شمسی ابتدا به میلادی تبدیل سپس به فرمت \lr{YYYY-MM-DD} فرمت‌بندی می‌شود.
\end{description}

\subsection{باگ ۴: HasIncompleteProfile (تغییر نام IsNewUser)}
\begin{description}[rightmargin=0.5cm]
  \item[مشکل:] فیلد \lr{IsNewUser} در پاسخ \lr{verify-otp} نام گمراه‌کننده‌ای داشت.
    این فیلد نشان می‌دهد که کاربر باید پروفایل خود را تکمیل کند، نه اینکه حتماً
    کاربر جدید است.
  \item[راه‌حل:] نام فیلد به \lr{HasIncompleteProfile} تغییر یافت.
    فرانت‌اند بر اساس این فیلد کاربر را به \lr{/complete-profile} هدایت می‌کند.
\end{description}

\newpage

%% ─── ۹. بهبود ورود ادمین ───────────────────────────────────────────────────
\section{بهبود امنیت ورود مدیر ارشد}

\subsection{مشکل}
صفحه ورود مدیر ارشد (\lr{/admin/login}) از همان فرم ورود با \lr{OTP}
استفاده می‌کرد که برای کاربران عادی طراحی شده بود.

\subsection{راه‌حل}
یک \lr{Endpoint} مستقل \lr{POST /api/admin/auth/login} با \lr{AdminLoginCommand}
ایجاد شد که:
\begin{itemize}[rightmargin=0.5cm]
  \item فقط با نام کاربری و رمز عبور کار می‌کند (نه \lr{OTP})
  \item رمز عبور با \lr{BCrypt} تأیید می‌شود
  \item توکن با نقش \lr{Admin} صادر می‌شود
  \item صفحه ورود ادمین در فرانت‌اند طراحی مستقل دارد
\end{itemize}

\newpage

%% ─── ۱۰. صفحات تکمیل‌شده ───────────────────────────────────────────────────
\section{صفحات تکمیل‌شده در این اسپرینت}

\subsection{صفحه تکمیل پروفایل (CompleteProfile)}
صفحه \lr{/complete-profile} برای کاربرانی که پس از ورود با \lr{OTP} پروفایل
ناقص دارند نمایش داده می‌شود:

\begin{itemize}[rightmargin=0.5cm]
  \item وارد کردن نام نمایشی (اجباری)
  \item آپلود تصویر پروفایل (اختیاری)
  \item این صفحه توسط \lr{ProtectedRoute} با نقش \lr{customer} محافظت می‌شود
  \item پس از تکمیل، به داشبورد مشتری هدایت می‌شود
\end{itemize}

\subsection{پروفایل مشتری (Profile)}
صفحه \lr{/customer-dashboard/profile}:
\begin{itemize}[rightmargin=0.5cm]
  \item نمایش اطلاعات فعلی کاربر
  \item ویرایش نام نمایشی
  \item تغییر تصویر پروفایل با آپلود
  \item تغییر رمز عبور (فرم جداگانه)
  \item نمایش شماره موبایل (فقط خواندنی)
\end{itemize}

\subsection{پروفایل مدیر رستوران (ManagerProfile)}
مشابه پروفایل مشتری با امکانات مشابه.

\subsection{مدیریت میز (TableManagement)}
صفحه مدیریت میز از اسپرینت قبل تکمیل شد:
\begin{itemize}[rightmargin=0.5cm]
  \item لیست میزهای فعلی با شماره و ظرفیت
  \item فرم افزودن میز جدید
  \item دکمه حذف با دیالوگ تأیید
  \item نمایش وضعیت اشغال بودن میز
\end{itemize}

\newpage

%% ─── ۱۱. تست نهایی ─────────────────────────────────────────────────────────
\section{تست نهایی و استقرار}

\subsection{تست‌های انجام‌شده}
\begin{longtable}{|r|r|r|}
\hline
\textbf{سناریو} & \textbf{نتیجه} & \textbf{توضیح} \\
\hline
\endhead
جریان کامل ثبت‌نام مشتری & موفق & \lr{OTP} + تکمیل پروفایل + داشبورد \\
\hline
جریان کامل ورود مشتری & موفق & \lr{OTP} + هدایت به داشبورد \\
\hline
ثبت رستوران توسط مدیر & موفق & فرم چندمرحله‌ای + ارسال برای تأیید \\
\hline
تأیید رستوران توسط ادمین & موفق & تغییر وضعیت به \lr{Approved} \\
\hline
رزرو کامل توسط مشتری & موفق & انتخاب تاریخ، میز، ارسال \\
\hline
تأیید رزرو توسط مدیر رستوران & موفق & تغییر وضعیت + اعلان \\
\hline
لغو رزرو توسط مشتری & موفق & تغییر وضعیت به \lr{Cancelled} \\
\hline
محدودیت نرخ \lr{OTP} & موفق & بعد از ۵ درخواست، \lr{429} \\
\hline
انقضا توکن و تجدید & موفق & \lr{Token Rotation} کار می‌کند \\
\hline
ورود ادمین & موفق & دسترسی به پنل ادمین \\
\hline
\end{longtable}

\subsection{محیط استقرار}
\begin{itemize}[rightmargin=0.5cm]
  \item بک‌اند: \lr{ASP.NET Core 8} روی \lr{Docker}
  \item پایگاه داده: \lr{PostgreSQL 16} روی \lr{Docker}
  \item کش: \lr{Redis 7} روی \lr{Docker}
  \item فرانت‌اند: \lr{Vite dev server} برای توسعه
  \item \lr{Docker Compose} برای اجرای همه سرویس‌ها با یک دستور
\end{itemize}

\newpage

%% ─── ۱۲. بازبینی ────────────────────────────────────────────────────────────
\section{بازبینی و گذر اسپرینت}

\subsection{Sprint Review}
\begin{itemize}[rightmargin=0.5cm]
  \item تمامی اهداف امنیتی محقق شدند
  \item تمام باگ‌های شناسایی‌شده رفع شدند
  \item سیستم آماده تحویل نهایی است
  \item جریان کامل سیستم (از ثبت‌نام تا رزرو و مدیریت) بدون خطا کار می‌کند
\end{itemize}

\subsection{Sprint Retrospective}

\textbf{چه چیزی خوب بود:}
\begin{itemize}[rightmargin=0.5cm]
  \item تمرکز بر کیفیت و امنیت به جای اضافه کردن ویژگی‌های جدید
  \item استفاده از استاندارد \lr{RFC 7807} فرمت پاسخ‌های خطا را یکدست کرد
  \item \lr{Token Rotation} امنیت سیستم را به طور قابل توجهی بهبود داد
\end{itemize}

\textbf{درس‌های آموخته‌شده از کل پروژه:}
\begin{itemize}[rightmargin=0.5cm]
  \item معماری تمیز از ابتدا هزینه توسعه را در درازمدت کاهش می‌دهد
  \item \lr{CQRS + MediatR} کدهای \lr{Business Logic} را از \lr{Controller} جدا و
    تمیز نگه می‌دارد
  \item استفاده از \lr{Redis} برای \lr{OTP} بسیار سریع‌تر و مناسب‌تر از پایگاه داده است
  \item \lr{Rate Limiting} باید از همان اسپرینت اول پیاده‌سازی شود
\end{itemize}

\subsection{آمار کلی پروژه}

\begin{longtable}{|r|r|}
\hline
\textbf{معیار} & \textbf{مقدار} \\
\hline
\endhead
تعداد \lr{Endpoint}های \lr{API} & بیش از ۳۵ \\
\hline
تعداد موجودیت‌های دامنه & ۹ موجودیت \\
\hline
تعداد \lr{Feature}های \lr{CQRS} & بیش از ۴۰ کامند/کوئری \\
\hline
تعداد صفحات فرانت‌اند & ۲۵ صفحه \\
\hline
تعداد کامپوننت‌های مشترک & بیش از ۱۵ کامپوننت \\
\hline
تعداد اسپرینت & ۳ اسپرینت \\
\hline
مدت کل پروژه & ۶ هفته \\
\hline
\end{longtable}

\end{document}
